% ==============================================================================
% Formation C# & SQL — Module 001 : Les Bases
% Fichier : src/P1_07_a.tex
% Sujet   : Chapitre 7 : Fonctions et méthodes
% ==============================================================================

\newpage
\section{Fonctions et Méthodes}
\marginpar{\tiny\color{gray}P1\_07\_a.tex}

En ingénierie logicielle, la création de méthodes (le terme C\# pour désigner les fonctions rattachées à une classe) est fondamentale pour garantir la modularité, la réutilisabilité et la lisibilité du code. Une méthode regroupe un ensemble d'instructions accomplissant une tâche algorithmique ou métier spécifique.

\subsection{Déclaration, Paramètres et Valeur de Retour}

La signature d'une méthode définit son comportement externe : son niveau d'accès, son type de retour, son identifiant et ses paramètres d'entrée. 
Si la méthode ne produit aucune donnée de sortie exploitable, son type de retour est déclaré comme \texttt{void}.

\begin{lstlisting}[style=csharpstyle]
// Méthode 'void' (aucune valeur retournée)
public void EcrireLog(string message)
{
    Console.WriteLine($"[LOG] {DateTime.UtcNow} : {message}");
}

// Méthode avec retour d'une donnée typée (decimal)
public decimal CalculerTTC(decimal montantHT, decimal tauxTVA)
{
    decimal montantTVA = montantHT * (tauxTVA / 100);
    return montantHT + montantTVA;
}
\end{lstlisting}

La clause \texttt{return} stoppe immédiatement l'exécution de la méthode et restitue la valeur à l'appelant. Le type de cette valeur doit correspondre strictement à la signature définie.

\subsection{Paramètres Optionnels et Surcharge}

C\# offre deux mécanismes distincts pour assouplir l'appel des méthodes :

\subsubsection{Valeurs par défaut (Paramètres Optionnels)}
Un paramètre peut se voir attribuer une valeur par défaut dans la signature. Il devient alors optionnel lors de l'appel. Les paramètres optionnels doivent obligatoirement être déclarés \textit{à la fin} de la liste des paramètres.

\begin{lstlisting}[style=csharpstyle]
public void ConfigurerConnexion(string ip, int port = 1433, int timeout = 30)
{
    Console.WriteLine($"Connexion à {ip}:{port} (Timeout: {timeout}s)");
}

// Appels valides :
ConfigurerConnexion("192.168.1.10"); // port = 1433, timeout = 30
ConfigurerConnexion("10.0.0.5", 8080); // timeout = 30
\end{lstlisting}

\subsubsection{Surcharge de Méthodes (\textit{Overloading})}
La surcharge consiste à définir plusieurs méthodes portant \textbf{le même nom}, mais ayant une \textbf{signature différente} (nombre, type ou ordre des paramètres). Le compilateur détermine la méthode à exécuter en analysant les arguments passés.

L'avantage principal de l'\textit{overloading} est de proposer une API naturelle. Au lieu d'inventer des noms comme \texttt{AfficherProfilBasique} et \texttt{AfficherProfilComplet}, on centralise l'action sous un seul verbe.

\begin{lstlisting}[style=csharpstyle]
// Version 1 : 1 paramètre (string)
public void AfficherProfil(string pseudo)
{
    Console.WriteLine($"Profil basique : {pseudo}");
}

// Version 2 : 2 paramètres (string, int) -> Le nombre change
public void AfficherProfil(string pseudo, int age)
{
    Console.WriteLine($"Profil complet : {pseudo}, {age} ans");
}

// Version 3 : 2 paramètres (int, string) -> L'ordre change (signature valide)
public void AfficherProfil(int age, string pseudo)
{
    Console.WriteLine($"[Mode inversé] : {pseudo}");
}
\end{lstlisting}

\subsection{Nombre Variable d'Arguments (\texttt{params})}

Le mot-clé \texttt{params} permet à une méthode d'accepter un nombre indéfini d'arguments du même type, qui seront encapsulés automatiquement dans un tableau. Ce mot-clé ne peut être appliqué qu'au tout dernier paramètre de la signature.

\begin{lstlisting}[style=csharpstyle]
public decimal CalculerMoyenne(params decimal[] valeurs)
{
    if (valeurs.Length == 0) return 0;
    
    decimal somme = 0;
    foreach (decimal val in valeurs)
    {
        somme += val;
    }
    return somme / valeurs.Length;
}

// L'appelant peut passer une liste d'arguments séparés par des virgules
decimal moyenne = CalculerMoyenne(12.5m, 15.0m, 18.5m);
\end{lstlisting}

\subsection{Passage de Paramètres : \texttt{ref} et \texttt{out}}

Par défaut, le passage d'arguments (pour les types valeur) s'effectue \textbf{par valeur} (copie isolée). Le modificateur \texttt{ref} exige un passage par référence (adresse mémoire), permettant à la méthode de muter directement la variable de l'appelant. La variable doit obligatoirement être initialisée avant l'appel.

Le modificateur \texttt{out} opère de la même manière, mais son rôle est d'exporter un résultat. Il exempte l'appelant d'initialiser la variable, mais force contractuellement la méthode à lui assigner une valeur avant l'instruction de retour.

\begin{lstlisting}[style=csharpstyle]
// Exemple d'utilisation du pattern TryParse avec 'out' inline
string fluxReseau = "1024";

if (int.TryParse(fluxReseau, out int portServeur))
{
    Console.WriteLine($"Port alloué : {portServeur}");
}
\end{lstlisting}
\textit{Note : Bien que \texttt{out} fut longtemps utilisé pour retourner de multiples valeurs, l'architecture C\# moderne préconise plutôt l'utilisation des \textbf{Tuples} pour cet usage.}

\subsection{Méthodes Statiques vs d'Instance}

\begin{itemize}
    \item \textbf{Méthodes d'instance} : Nécessitent l'instanciation préalable de la classe (la création d'un objet en mémoire). Elles peuvent manipuler l'état interne (les champs) propre à cette instance spécifique.
    \item \textbf{Méthodes statiques (\texttt{static})} : Sont attachées au type lui-même (la classe) et non à une instance. Elles ne peuvent pas accéder aux données d'une instance. Elles sont conçues pour des opérations utilitaires pures (ex: \texttt{Math.Abs()} ou \texttt{int.TryParse()}).
\end{itemize}

\begin{tcolorbox}[colback=backcolour,colframe=keywordcolor,title=\textbf{Pièges Courants \& Erreurs}]
\begin{itemize}
    \item \textbf{Ambiguïté de surcharge} : Créer deux méthodes avec des signatures identiques mais des types de retour différents génère une erreur de compilation, car le type de retour ne fait pas partie de la signature permettant de résoudre la surcharge.
    \item \textbf{Rupture du contrat \texttt{out}} : \texttt{CS0177} survient inévitablement si la méthode omet d'affecter une valeur au paramètre \texttt{out} avant de se terminer.
    \item \textbf{Ordre des paramètres optionnels} : Tenter de définir un paramètre optionnel \textit{avant} un paramètre requis provoquera une erreur stricte du compilateur.
\end{itemize}
\end{tcolorbox}

\subsection{Synthèse}
\begin{itemize}
    \item \textbf{\texttt{void} / Type de retour} : Définit si la méthode opère comme une action silencieuse ou comme une fonction calculant un résultat.
    \item \textbf{Surcharge \& Paramètres optionnels} : Flexibilisent l'API de votre classe et facilitent l'appel de vos méthodes.
    \item \textbf{\texttt{params}} : Simplifie le passage d'une liste dynamique de variables du même type.
    \item \textbf{\texttt{ref} / \texttt{out}} : Permettent le passage par référence, utiles pour muter la variable appelante ou implémenter le pattern \texttt{TryParse}.
    \item \textbf{\texttt{static}} : Rend la méthode globale et indépendante de tout état d'instance.
\end{itemize}
